\chapter{Discussion}
\label{chap:discussion}

\section{Interpretation of Findings}

The results separate four kinds of evidence: whether the SAE preserves the
activation, whether its features predict the decision, what those features
respond to, and whether editing them changes the decision. The SAE features
pass the second and fourth tests only partly, and the third test shows that
they are not simple ``deference features''.

\paragraph{The decision is readable after the feedback, not before.} A probe
on layer~25 predicts whether the Solver will adopt the Critic's answer with
0.843 balanced accuracy, well above probes on the visible text (0.623), on
metadata such as candidate order and data source (0.525), and on the Solver's
state before it saw the feedback (0.597). The decision therefore depends on
how the model processes the feedback, not only on surface features of the
prompt. The SAE bases keep most but not all of this signal: 0.773 for the
custom $L_1$ SAE and 0.706 for the pretrained family
(Table~\ref{tab:probe-performance}). Some of what a dense probe reads is lost
when the activation is rewritten as sparse features.

\paragraph{The selected features encode the inputs to the decision more than
the decision itself.} For only 4 of the 80 selected SAE features was ``the
Solver adopted'' the episode property they tracked most strongly. Most
tracked the kind of answer the Critic gave: a date for 11 of the 20 pretrained
features, and an objection that the question cannot be answered for 9 of the
20 $L_1$ features (Section~\ref{sec:res-predictive}). Their correlation with
the decision itself was about $\pm0.1$. A plausible reading is that the Solver
decides from properties of the two candidates, and the probe found the
features that represent those properties. The features are relevant to the
decision without being a description of it.

\paragraph{The causal effect is real but small, and it works by tipping close
decisions.} In the prespecified tests, targeted edits of the pretrained and
decoder-adapted SAE features flipped 1.6 to 2.5 percentage points more
decisions than random edits of the same size, and four of the nine tests
survived Holm correction (Table~\ref{tab:primary-causal}). The flips were
fluent switches to the other candidate, not broken text, and they came mainly
from episodes where the Solver was less certain of its answer. Effects were
stronger when pushing the Solver away from adopting than toward it. We read
this as evidence that the selected directions take part in the decision, but
as one small contribution among many: the decision is not controlled by ten
features.

\paragraph{Interpretability and causal leverage pulled in opposite
directions.} The custom $L_1$ features were the easiest to describe (median
validation correlation of Claude's descriptions 0.77, against 0.38 for the
pretrained features), yet they were the weakest steering handles. Part of the
reason is geometric: the $L_1$ Adoption and Retention directions are closely
aligned (cosine 0.754), so they largely encode the same thing with opposite
signs, and the arm that combines them cancels. The pretrained features were
harder to name but steered best. A feature that is easy to explain from its
activating examples is therefore not necessarily one through which the model
acts.

\paragraph{The raw basis steers harder; the SAE basis edits more cleanly.}
Equal-size edits along probe-weighted raw neurons gave larger point effects,
up to $+4.1$ points, but the strongest raw edit left 10.1\% of answers outside
both candidates. Only one of nine post hoc comparisons found raw significantly
ahead of an SAE. The SAE features buy selectivity, not strength.

\paragraph{The LLM-edited SAE mostly shows the strength of contrastive
directions.} The exploratory LLM-edited SAE produced the largest effects of
any representation ($+8.5$ to $+9.6$ points). Its own controls explain most
of this. A mechanical edit that rotates the same 20 decoder rows toward the
train mean difference between Adoption and Retention reached $+6.2$ to $+6.6$
points, while rotating them by the same angles in random directions reached
only $+0.6$ to $+1.4$.
% Source: results/intervention_llm_edited_layer25/20261112T005617641082Z-f48aae4eb4fa/merge/table_P_primary.md
The large effect comes mostly from moving the features toward a contrastive
steering direction, and it raised Out answers to 7--12\%
(Section~\ref{sec:res-causal}). It shows that the decision can be steered
along a learned direction; it does not show that the edited features have a
clearer meaning.

\paragraph{Probe-based checks overstate causal effects.} The earlier designs
(Section~\ref{sec:res-exploratory}) showed the same pattern repeatedly: an
edit of 1\% of the activation norm moved the probe's prediction by 2.5 points
but left generated answers unchanged, and probe-scored flips appeared even
when no feature was edited. The offline ablation in
Section~\ref{sec:res-predictive} is useful for choosing features, but only
generation measures what the model does.

\section{Comparison with Prior Work}

Our results agree with recent evaluations that set SAEs against strong
baselines. Probes on SAE features fell short of probes on raw activations, as
in the sparse-probing case study of Kantamneni et al.\
\cite{kantamneni2025useful}, and raw directions steered at least as strongly
as SAE features, in line with AxBench \cite{wu2025axbench}. Korznikov et al.\
found that random baselines can approach trained SAEs on several measures
\cite{korznikov2026sanity}. In our study, random-feature edits of the same
size accounted for most of the flips that targeted edits produced, which is
why every effect is reported as targeted minus random.

Arad et al.\ find that SAE steering works much better when features are
selected for their effect on the output rather than for the inputs they
respond to
\cite{arad2025steering}. Our selection used probe weights, an input-side
criterion, and it produced features that were describable but steered weakly.
The trade-off we observed between description quality and steering strength
is consistent with their argument. Feature-absorption studies make a related
point: a feature can look coherent in its top examples and still fail to track
the concept consistently \cite{chanin2024absorption}. Scoring each description
on held-out episodes, as we did, guards against this but does not remove it.

On the behavioral side, our pipeline reproduced known patterns of model
feedback. When Gemma reviewed its own answer, it repeated that answer in 179
of 183 reviews, consistent with the limits of intrinsic self-correction
\cite{huang2024selfcorrect}. Under free-form feedback the Solver adopted the
Critic's answer in 91.7\% of disagreements, and under the candidate format it
adopted wrong corrections more often than right ones
(Section~\ref{sec:problems}), echoing sycophancy findings in which models
follow an expressed view at the cost of accuracy \cite{sharma2024sycophancy}.
What this study adds is an internal view: the uptake decision is linearly
readable at layer~25, but the readable signal is mostly about the form of the
Critic's answer. Finally, the strength of the mechanical edit, which moves
features toward a mean-difference direction, matches the success of
contrastive activation addition \cite{panickssery2023caa}.

\section{Threats to Validity}

\paragraph{The LLM judge.} All behavior labels, including the labels of edited
generations, come from Claude Sonnet~4.5. On 100 audited validation episodes,
human and Claude labels agreed on 84 ($\kappa = 0.787$), and no episode was
labeled Adoption on one side and Retention on the other
(Section~\ref{sec:human-audit}). The audit did not cover edited generations,
which may be less typical than the stored answers. Two facts limit this risk:
most flips were fluent choices of the other candidate, and Out answers, the
class Claude labels least precisely, are reported separately and never count
as flips.

\paragraph{The synthetic data shift the distribution.} Synthetic questions
supply 61\% of the binary episodes but produce a disagreement far less often
than MuSiQue questions (27.4\% against 62.7\%), and their hop mix differs
(Section~\ref{sec:synthetic-data}). The probes therefore learn mostly from
synthetic examples. In a post hoc split, the causal effects appeared on
MuSiQue episodes (three tests surviving Holm, $n = 371$) and not on synthetic
episodes (none, $n = 612$). This is reassuring for the natural questions but
means the synthetic data added volume to the probes without adding causal
evidence. We cannot tell whether synthetic decisions are simply more certain
or whether the features represent the two sources differently; the split was
not prespecified.

\paragraph{Candidate order and source.} Listing the Critic's answer first
raised Adoption by 6.4 points on MuSiQue and lowered it by 8.1 points on
synthetic questions (Table~\ref{tab:order-final}). On their own, order and
source predict the decision poorly (metadata probe 0.525), but the two
strongest $L_1$ Adoption features were tied most closely to which candidate
was listed first. Some selected features therefore encode presentation rather
than uptake, and the causal effects have not been broken down by order.

\paragraph{Broken outputs.} An edit can change the label by degrading the
answer instead of changing the choice. We separated switches from Out answers,
and restricting the analysis to answers that chose a candidate left the SAE
effects essentially unchanged. The raw-activation control and the LLM-edited
SAE did raise Out rates, so their larger effects are partly of this kind.
Seventeen episodes were excluded because a random edit produced an empty
answer that could not be labeled.

\paragraph{Edit size.} All primary edits used one size, 20\% of the activation
norm, chosen on a 50-episode pilot by a prespecified rule. The earlier clamped
design showed how unequal edit sizes can manufacture a 73-point effect, and
equal-size matching removes that confound. With a single size, however, we do
not know how the effect grows with the edit.

\paragraph{Deference or discrimination.} The study cannot fully separate a
Solver that defers to any recommendation from one that recognizes good
feedback (Section~\ref{sec:problems}). The controlled conditions that would
have fixed whether the Critic was right were dropped, and natural feedback
leaves the Critic's correctness confounded with question difficulty. The
observation that the Solver adopted wrong corrections more often than right
ones argues against discrimination, but it is an association, not a
controlled test, and it was measured with the rule-based labels.

\paragraph{Multiple comparisons.} The primary claim rests on nine prespecified
tests with Holm correction. The results by direction, by source and for the
LLM-edited SAE are secondary or exploratory, and the $L_1$ \textsc{decrease}
arm narrowly missed correction ($p_{\mathrm{Holm}} = 0.051$). We do not count
these as confirmed effects.

\section{Limitations}

Section~\ref{sec:model-limitations} lists the limits that come from the
models: their size, the single layer and token position, the imperfect SAE
reconstruction, limited compute, and greedy decoding. The study design adds
four more.

\begin{itemize}
    \item \textbf{A binary target.} The analysis explains the Solver's choice
    only when the two answers differ. Agreement episodes, about 64\% of the
    collection (Table~\ref{tab:final-labels}), and Out answers fall outside
    the target.
    \item \textbf{One feature-selection rule.} Each set holds the ten features
    ranked highest by mean absolute SHAP value on the probe. Other rules,
    especially ones that rank features by their effect on generation, could
    find stronger handles.
    \item \textbf{One intervention site and one edit size.} Edits were applied
    once, at the final prompt token, at a single relative norm.
    \item \textbf{One task family and one hosted judge.} All questions are
    short-answer multi-hop questions in the MuSiQue format, and regenerating
    the labels requires the same pinned version of a hosted model that the
    team does not control.
\end{itemize}

\section{Implications for Other Researchers}

\begin{itemize}
    \item \textbf{Measure effects in generation.} Probe-scored ablations are
    useful for choosing features but can overstate causal effects by an order
    of magnitude or more.
    \item \textbf{Use matched controls from the start.} Reconstruction-only
    passes, equal-size random edits and a separate count of broken answers
    each exposed an apparent effect in our earlier designs that was not a
    changed decision. The clamped design alone reported a 73-point effect.
    \item \textbf{Randomize presentation and report by source.} Candidate
    order shifted adoption by several points in opposite directions in the two
    sources, and pooling hid it.
    \item \textbf{Do not equate interpretability with control.} The most
    describable features steered least; both properties need their own test.
    \item \textbf{Audit the labels and pilot the behavior.} A frozen judge
    with recorded settings and a blind human audit made the LLM-derived labels
    usable. Small pilots caught protocols that would have wasted a full
    collection: free-form feedback produced almost no Retention, and a Critic
    reviewing its own model's answer copied it.
\end{itemize}

\section{Next Steps}

The following extensions would address the main open questions. None has
been run.

\begin{itemize}
    \item Restore controlled conditions in which the Critic is known to be
    right or wrong, to separate deference from discrimination.
    \item Select features by their measured effect on generation rather than
    by probe weight, following Arad et al.\ \cite{arad2025steering}.
    \item Edit at every generated token and at more than one layer, and test
    several edit sizes to measure a dose--response curve.
    \item Prespecify the analysis by data source and candidate order, with a
    larger MuSiQue intervention set.
    \item Add a contrastive mean-difference direction as a raw steering
    baseline, since the mechanical edit suggests it is a strong comparator.
    \item Repeat the study with larger Solvers and other model families.
\end{itemize}
